% !TeX program = xelatex
% Расписка студента (NDA) — официальная форма НИУ ВШЭ (Расписка студента NDA 2026).
% Предоставляется на защиту ВКР, если результаты носят конфиденциальный характер.
% Персональный документ: у каждого автора своя расписка (ФИО и тема ВКР
% подставляются из \hseAuthorName / \hseProjectName пер-автор). Поля паспорта,
% адреса, компании и подписи заполняются от руки после печати.
\documentclass[12pt,a4paper]{extarticle}
\input{../common/preamble}
\newcommand{\docname}{Расписка}
\newcommand{\doccode}{}
\newcommand{\doccodelu}{}
\input{../common/commands}
\begin{document}
\pagestyle{empty}
\sloppy
\setlength{\parindent}{0pt}
\setlength{\parskip}{6pt}

\begin{center}
  \textbf{\large Расписка}
\end{center}

\vspace{0.5em}

Я, \hseAuthorName~\hrulefill

\vspace{0.3em}
паспорт \hrulefill

\vspace{0.3em}
выдан \hrulefill

\vspace{0.3em}
проживающий по адресу \hrulefill

\hrulefill

\vspace{0.5em}
в связи с предоставлением мне \underline{\hspace{5cm}} (далее~--- Компания)
доступа к~определённой конфиденциальной информации добровольно обязуюсь:

\begin{enumerate}[leftmargin=1.4em, itemsep=5pt]
  \item Не совмещать работу в~НИУ ВШЭ над выпускной квалификационной работой
        «\hseProjectName» (далее~--- ВКР) с~работой в~иных организациях (кроме
        Компании и~НИУ ВШЭ).

  \item Не использовать в~рамках работы над ВКР, не~включать в~состав
        результатов ВКР какие-либо сторонние (в~том числе свои, созданные до
        начала ВКР) исходные коды программ, картинки, тексты, on-line сервисы и~любые
        другие аналогичные материалы, которые могут являться объектом
        исключительного права, без специального согласования с~руководителем /
        консультантом ВКР от~Компании.

  \item Не передавать другим организациям (включая НИУ ВШЭ, не~включая Компанию)
        и~частным лицам, как во~время работы над ВКР, так и~после окончания
        работы, любые полученные в~рамках выполнения этой ВКР материалы, что
        включает в~себя: блок-схемы и~диаграммы с~архитектурными решениями,
        собственно программные коды, описания API, документацию для разработчиков
        и~администраторов, перечни компонентов и~зависимостей, сведения о~применяемых
        алгоритмах для защиты информации, сведения из~системы отслеживания ошибок
        (bug tracking), результаты нагрузочного и~сценарного тестирования,
        относящиеся как к~подлежащему разработке программному обеспечению, так и~к~уже
        существующему программному обеспечению, в~том числе к~существующим
        коммерческим продуктам и~онлайн-сервисам Компании, независимо от~наличия на~них
        грифа «Конфиденциально» и~подобных (далее~--- Материалы), а~также любую
        ставшую мне известной информацию касательно внутренней деятельности Компании
        (включая порядок организации работы, порядок и~размеры оплаты труда,
        организационную структуру фирмы и~пр.).

  \item Не использовать созданные самостоятельно в~рамках выполнения ВКР Материалы
        в~других разработках за~пределами НИУ ВШЭ. Не~передавать на~сторону и~не
        использовать для других целей любые программные продукты, SDK, библиотеки,
        инструментальные и~отладочные средства, полученные и/или используемые в~процессе
        выполнения ВКР, за~исключением свободного программного обеспечения (FOSS),
        обнародованного и~свободно распространяемого третьими лицами под одной из
        открытых лицензий, список которых приведён на~сайте
        \url{https://spdx.org/licenses}.

  \item Не публиковать программные коды, разработанные в~рамках выполнения ВКР;
        публиковать другие разработанные в~рамках выполнения ВКР Материалы, в~том
        числе в~составе научных публикаций, и~демонстрировать их~в~рамках защиты
        данной ВКР и~иных проектов только после согласования с~руководителем или
        консультантом ВКР от~Компании.

  \item Вышеизложенные обязательства принимаются мною на~срок всего обучения в~НИУ
        ВШЭ, а~обязательства по~неразглашению Материалов, изложенные в~пунктах
        3--5,~--- на~срок всего обучения в~НИУ ВШЭ и~плюс 3~(три) года после
        прекращения обучения в~НИУ ВШЭ.

  \item Участвовать в~защите ВКР в~НИУ ВШЭ только в~присутствии руководителя или
        консультанта ВКР от~Компании.
\end{enumerate}

\vspace{1.5em}

\noindent
Дата: «\underline{\hspace{1.2cm}}»~\underline{\hspace{3.5cm}}~20\underline{\hspace{0.6cm}}~г.
\hfill
Подпись: \underline{\hspace{4.5cm}}

\end{document}
